\subsection{RGSW}
The techniques above are used in the \gls{gsw} scheme~\cite{cryptoeprint:2013/340} and its ring variants~\cite{cryptoeprint:2014/816, cryptoeprint:2018/421}, where the ciphertext itself is built around a gadget. We now formally define an \gls{rgsw} ciphertext, and the operations required to employ \gls{rgsw} in \gls{fhe} schemes. %% TODO: Rewrite for style + TFHE approximates gadget + Formalize sample vs ciphertext vs encryption.

\begin{definition}[Gadget Matrix]
    \label{def:gadget-matrix}
    A gadget $g \in R_Q^\ell$ is extended to $k$-dimensional ciphertexts by the block-diagonal gadget matrix $G = I_k \otimes g \in R_Q^{k\ell\times k}$
\end{definition}

\begin{definition}[RGSW Ciphertext]
    Given a gadget matrix $G = I_2 \otimes g$, we define the \emph{masking} matrix $Z \in R_Q^{2\ell\times 2}$ as a collection of $2\ell$ \gls{rlwe} public keys $\Call{RLWE}{0}$. An \emph{\gls{rgsw}} encryption $C$ of $\mu \in R_Q$ is denoted $\Call{RGSW}{\mu}$ and defined as
    \begin{equation}
        \label{def:rgsw}
        C = Z + \mu\cdot G % \in R_Q^{2\ell\times 2}
    \end{equation}
\end{definition}

% Each row of an \gls{rgsw} ciphertext is a valid \gls{rlwe} sample. The upper block of $\ell$ rows are equivalent to $\Call{RLWE}{\mu g_i}$ for row $i$. The lower block of the remaining $\ell$ rows encode the message in the public components, where the $i$-th row effectively becoming $\Call{RLWE}{-\mu g_i\cdot s}$.

An \gls{rgsw} ciphertext encodes the gadget-scaled messages $\mu g_i$ row by row. In the top block, $\mu g_i$ is added to the masked element, replacing the zero message of the underlying $\Call{RLWE}{0}$; in the bottom block it is added to the public element, so the row decrypts to $-\mu g_i\cdot s$ instead. Every row of $C$ is therefore itself an \gls{rlwe} ciphertext.

Two \gls{rgsw} ciphertexts add homomorphically, just as \gls{rlwe} ciphertexts do: the zero encryptions and the gadget terms add component-wise, so $C_A + C_B$ encrypts $\mu_A + \mu_B$. Addition alone is, however, not sufficient for \gls{fhe}. This motivates the need for an \gls{rgsw}$\times$\gls{rgsw} product and a \gls{rlwe}$\times$\gls{rgsw} product. These are called the internal and external products, respectively, in \cite{cryptoeprint:2018/421}. First, we extend the gadget property to two dimensions:
% The multiplication mirrors the vector--matrix relationship between the two ciphertext types: the \gls{rlwe} ciphertext $x = (x_0, x_1)$ is mapped by a two-dimensional decomposition $G^{-1}(x)$ to a row vector of $2\ell$ small ring elements, which multiplies the $2\ell\times 2$ matrix $C$. The only requirement on $G^{-1}(x)$ is that it approximately inverts the $G$ matrix: 
\begin{equation}
    \label{eq:gadget-matrix-property}
    G^{-1}(x)\cdot G = x + e \pmod{Q}%, \qquad\|G^{-1}(x)\|_{\infty} \leq \beta
\end{equation}

%% TODO: Rewrite if time
No particular decomposition is prescribed. The canonical choice is the base-$B$ digit decomposition of the ciphertext coefficients~\cite{cryptoeprint:2018/421}, and it is the one usually meant when the external product is written down, but nothing in the definition depends on it. Any $G^{-1}(x)$ meeting \eqref{eq:gadget-matrix-property} at quality $\beta$ yields a valid product, and the noise bounds below are stated in terms of $\beta$ alone.

\begin{definition}[External Product]
    \label{def:prod:external}
    Let $x \in R_Q^2$ be an \gls{rlwe} encryption of $m \in R_t$, let $C \in R_Q^{2\ell\times 2}$ be an \gls{rgsw} encryption of $\mu \in R_Q$ under the same secret key, and let $G^{-1}$ be a gadget decomposition satisfying \eqref{eq:gadget-matrix-property}. The {external product} is defined as
    \begin{align}
        \boxdot: \Call{RLWE}{} \times \Call{RGSW}{} &\longmapsto \Call{RLWE}{} \nonumber \\
        \label{def:extprod}
        (x, C) &\longmapsto x \boxdot C = G^{-1}(x) \cdot C,
    \end{align}
    which evaluates $\Call{RLWE}{m}\,\boxdot\,\Call{RGSW}{\mu} = \Call{RLWE}{m\cdot\mu}$. Formal correctness and noise bounds are given in~\cite{cryptoeprint:2018/421}, though informally this follows from inserting \eqref{def:rgsw} into \eqref{def:extprod} and using the gadget property \eqref{eq:gadget-matrix-property}, yielding a sum of two \gls{rlwe} samples:
    \begin{equation}
        \label{eq:extprod-expansion}
        G^{-1}(x)\cdot C = G^{-1}(x)\cdot Z + \mu\cdot(x + e)
    \end{equation}
\end{definition}

\vspace{1pt}

\begin{lemma}[Gadget Matrix Construction]
    \label{lemma:gadget-matrix-construction}
    Two vectors $P(u), D(v)$ that satisfy the gadget property \eqref{def:gadget-property} define a valid pair of gadget matrices \eqref{def:Gmat} and \eqref{def:Ginvmat} satisfying \eqref{eq:gadget-matrix-property}, with $\|G^{-1}(x)\|_\infty \leq \beta$ inherited from the quality of $D(v)$.
    \begin{gather}
        \label{def:Gmat}
        G = I_2 \otimes P(1) \\
        \label{def:Ginvmat}
        G^{-1}(x) = [D(x_0)|D(x_1)]
    \end{gather}

    \begin{proof}
        The block structure reduces the product to one inner product per component, each an instance of \eqref{def:gadget-property}:
        \begin{gather*}
            G^{-1}(x)\cdot G = [D(x_0)|D(x_1)]
            \begin{pmatrix}
                P(1) & 0 \\
                0 & P(1)
            \end{pmatrix}, \quad \|D(\cdot)\|_\infty \leq \beta \\
            \big(\langle D(x_0), P(1)\rangle,\; \langle D(x_1), P(1)\rangle\big)
            = (x_0 + e_0,\; x_1 + e_1) = x + e \pmod{Q}
        \end{gather*}
    \end{proof}
\end{lemma}

\begin{remark}%[Asymmetric Noise Growth]
    \label{rem:extprod-asymmetry} %% TODO: Include more sources/authors
    \gls{rgsw} ciphertexts have been by several authors, including \cite{cryptoeprint:2013/541, cryptoeprint:2018/421}, to be uniquely suited to evaluate long chains of products, due to an asymmetry in the noise propagation --- an asymmetry explained in greater detail in \cref{appendix:asymmetry}. Essentially, since $\|G^{-1}(x)\|_\infty \leq \beta$ and $\mu\in\{0,1\}$ we see that \eqref{eq:extprod-expansion} is bounded by $\beta\cdot Z + x + e$. The total error in the external product becomes dominated by the product $\beta \cdot Z$, while $x$ is effectively ignored. A chain of external products $(x\boxdot C_0) \boxdot C_1 \boxdot ...$ is thus able to sustain itself with minimal noise growth as long as the \gls{rgsw} ciphertexts $C_i$ are fresh encryptions.
\end{remark}

\begin{definition}[Internal Product]
    \label{def:prod:internal}
    Let $A, C \in R_Q^{2\ell\times 2}$ be \gls{rgsw} encryptions of $\mu_A, \mu_C \in R_Q$ under the same secret key and gadget, and let $A_i$ denote the $i$-th row of $A$. The \emph{internal product} applies the external product to each row of the left operand:
    \begin{align}
        \boxtimes: \Call{RGSW}{} \times \Call{RGSW}{} &\longrightarrow \Call{RGSW}{} \nonumber \\
        (A, C) &\longmapsto A \boxtimes C = \begin{bmatrix}
            A_0 \boxdot C \\
            A_1 \boxdot C \\
            \vdots \\
            A_{2\ell - 1} \boxdot C
        \end{bmatrix},
    \end{align}
    which returns an encryption of the product of the messages: $\Call{RGSW}{\mu_A} \boxtimes \Call{RGSW}{\mu_C} = \Call{RGSW}{\mu_A\cdot\mu_C}$.
\end{definition}

One internal product amounts to $2\ell$ external products, and correctness follows row by row: each external product multiplies the content of $A_i$ by $\mu_C$, so the stacked result is again of the form \eqref{def:rgsw} with message $\mu_A\cdot\mu_C$.

The asymmetry of \autoref{rem:extprod-asymmetry} applies to each of the $2\ell$ rows, and dictates the order in which a chain of products should be evaluated. The noise of the right-hand operand is amplified by the decomposition in every row of the result, while the noise of the left-hand operand is only carried through. When accumulating a running product, the accumulator must therefore be kept on the left with each fresh ciphertext supplied on the right, so that the noise grows additively rather than multiplicatively.

\begin{remark}[Operand Order]
    \label{rem:operand-order}
    The convention above places the \gls{rlwe} operand of the external product on the left and the \gls{rgsw} operand on the right, so that both products read left to right with the accumulator first. TFHE writes the external product the other way round~\cite{cryptoeprint:2018/421}, with the \gls{rgsw} ciphertext on the left. The operation is the same, but any expression carried over from that literature must be flipped before it is read against \eqref{def:extprod}.
\end{remark}
